\documentclass[11pt]{article}
\usepackage{amsmath,amssymb,amsthm,hyperref}
\newtheorem{theorem}{Theorem}
\newtheorem{lemma}{Lemma}
\title{Excluding scaled monic quadratic numerator profiles in Erd\H{o}s 243}
\author{Research notes; no novelty claim}
\date{6 November 2026}
\begin{document}
\maketitle

\noindent\textbf{Status.} This is a restricted exclusion, not a proof of
Problem 243. It does not assert that the numerator of an arbitrary rational
pseudo-greedy tail is polynomial. The integer-tail framework is from
\href{https://math.colgate.edu/~integers/aa28/aa28.pdf}{Koizumi (2026)}.
Gcd stabilization, adjacent coprimality, and the
three-consecutive-numerator square obstruction are used to restrict
polynomial profiles in
\href{https://github.com/wcook04/plectis-erdos/blob/main/paper/243/erdos243-reciprocal-tail-reasoning-surface.tex}{Cook's companion paper}.
No claim of priority is made here.

\begin{theorem}
There are no positive integer sequences $C_n,D_n,a_n$ satisfying
\[
 C_{n+1}=a_n C_n-D_n>0,\qquad D_{n+1}=a_nD_n
\]
for which, eventually,
\[
 C_n=t(n^2+bn+c),\qquad t\in\mathbb Q_{>0},\quad b,c\in\mathbb Z.
\]
In particular, this entire class of exact numerator profiles cannot
produce a counterexample to Erd\H{o}s 243.
\end{theorem}

\begin{proof}
On the asserted tail, put $P(n)=n^2+bn+c$ and $F_n=D_n/t$.
The identity $F_n=a_nP(n)-P(n+1)$ shows that $F_n$ is a positive
integer. Thus we may work with the integer states $P(n),F_n$ and retain
the same multipliers. Write $G_n=\gcd(P(n),F_n)$.
The recurrences imply $G_n\mid G_{n+1}$, and $G_n$ divides every
subsequent numerator $P(j)$. Therefore $G_n$ divides the second
difference $P(n+2)-2P(n+1)+P(n)=2$. The divisibility chain stabilizes
at a value $g\in\{1,2\}$. On a further tail put
\[
 U_n=P(n)/g,\qquad V_n=F_n/g.
\]
These are positive integers, satisfy the same two recurrences, and obey
\[
 \gcd(U_n,V_n)=1,\qquad\gcd(U_n,U_{n+1})=1.
\]
The latter follows from $V_n=a_nU_n-U_{n+1}$.

Suppose first that $b$ is even. The parities of $P(n)$ alternate, so
$g=1$. If $2$ divided any sufficiently late $F_n$, it would divide
all subsequent $F_j$, contradicting primitive coprimality at the
infinitely many even $P(j)$. Thus all sufficiently late $F_n$ are odd,
as are $a_n$ because $F_{n+1}=a_nF_n$. Eliminating $F_n$ gives
\begin{equation}\label{eq:three}
 U_{n+2}=(a_n+a_{n+1})U_{n+1}-a_n^2U_n.
\end{equation}
Choose a sufficiently late $n$ with $P(n)$ odd. Then $P(n+1)$ is even,
so the first term on the right is divisible by $4$, and the second is
congruent to $-P(n)$ modulo $4$. But
$P(n+2)-P(n)=4n+4+2b$ is divisible by $4$. This would give
$2P(n)\equiv0\pmod4$, a contradiction.

We may therefore take $b=2h+1$. Shifting the integer index by $h$
puts the polynomial into the form $n^2+n+d$, where
$d=c-h(h+1)\in\mathbb Z$. If $d$ is odd then $g=1$. If $d$ is even,
all numerator values are even, so adjacent primitive coprimality forces
$g=2$. Consequently $d/g$ is an integer. It must equal $1$ or $-1$:
otherwise choose a prime $p$ dividing $d/g$, with any prime allowed if
$d=0$, and choose arbitrarily late $n$ divisible by $gp$. Both
\[
 U_n=\frac{n(n+1)}g+\frac dg,\qquad
 U_{n-1}=\frac{(n-1)n}g+\frac dg
\]
would then be divisible by $p$, a contradiction.

There remain just the following four profiles:
\[
 Q_1(n)=n^2+n+1,\quad Q_2(n)=n^2+n-1,\quad
 Q_3(n)=\frac{n^2+n+2}{2},\quad
 Q_4(n)=\frac{n^2+n-2}{2}.
\]
At an index with $U_n\equiv0\pmod p$, equation~\eqref{eq:three}
implies
\begin{equation}\label{eq:square}
 -U_{n-1}U_{n+1}\equiv(a_{n-1}U_{n-1})^2\pmod p.
\end{equation}
Three profiles are immediately excluded by this square condition:
\[
\begin{array}{c|c|c|c|c}
 \text{profile}&p&n\bmod p& (U_{n-1},U_{n+1})\bmod p
 &-U_{n-1}U_{n+1}\bmod p\\\hline
 Q_1&3&1&(1,1)&2\\
 Q_3&7&3&(4,4)&5\\
 Q_4&3&1&(2,2)&2
\end{array}
\]
The last entries are nonsquares in the indicated fields. Each residue
class contains arbitrarily late indices, so the eventual nature of the
profile causes no difficulty. For $Q_3,Q_4$, division by $2$ is valid
modulo these odd primes.

Finally take $U_n=Q_2(n)$ and choose a sufficiently late $n\equiv7\pmod{11}$.
The four consecutive values at $n-2,n-1,n,n+1$ reduce to
\[
 (7,8,0,5)\pmod{11}.
\]
The recurrence at $n$ gives $V_n\equiv-5\equiv6$. The exact identity
\begin{equation}\label{eq:Didentity}
 U_j V_{j+1}=V_j(V_j+U_{j+1})
\end{equation}
follows by substituting $a_jU_j=V_j+U_{j+1}$ into
$V_{j+1}=a_jV_j$. At $j=n-1$ it implies
$V_{n-1}^2\equiv8\cdot6\equiv4\pmod{11}$, hence
$V_{n-1}\equiv2$ or $9$. At $j=n-2$ it instead requires a solution $x=V_{n-2}$ of
\[
 x^2+8x-7V_{n-1}\equiv0\pmod{11}.
\]
Its discriminant is $64+28V_{n-1}\equiv10$ or $8$ respectively.
Neither is in the list $\{0,1,3,4,5,9\}$ of squares modulo $11$.
This excludes the last profile and proves the theorem.
\end{proof}

\paragraph{A further class, with arbitrary rational scale and offset.}
We use the following standard consequence of Chebotarev. This is the
square-specialisation lemma used in Cook's companion paper, stated here
only in the form needed below.

\begin{lemma}\label{lem:specialise}
Let $f\in\mathbb Q[X]$ have an algebraic root $\alpha$, and let
$H\in\mathbb Q[X]$ with $H(\alpha)\ne0$. Suppose that for all but
finitely many rational primes $p$, every root $r\in\mathbb F_p$ of the
reduction of $f$ has $H(r)$ a square. Then $H(\alpha)$ is a square
in $\mathbb Q(\alpha)$.
\end{lemma}
\begin{proof}
Put $K=\mathbb Q(\alpha)$. If $H(\alpha)$ were not a square, adjoining
$\beta=\sqrt{H(\alpha)}$ would give a quadratic extension of $K$.
In a finite Galois extension $L/\mathbb Q$ containing it, an automorphism
fixing $K$ and sending $\beta$ to $-\beta$ exists by extension of
automorphisms. Chebotarev gives infinitely many rational primes with
Frobenius conjugate to that automorphism. Choose a prime of $L$ above
each of them having that exact Frobenius. After removing the finitely
many denominator, ramification, characteristic-two, and zero-reduction
exceptions, $\overline\alpha$ is fixed by Frobenius and lies in
$\mathbb F_p$, whereas neither $\overline\beta$ nor
$-\overline\beta$ is fixed. This contradicts the asserted square
condition at the root $\overline\alpha$.
The Chebotarev input is the infinitude statement for each conjugacy
class; see, for example,
\href{https://www.jmilne.org/math/CourseNotes/CFT.pdf}{Milne,
Class Field Theory, Chapter V, Theorem 3.23}.
\end{proof}

\begin{theorem}\label{thm:positiveoffset}
No positive integer orbit as in the first theorem can have, eventually,
\[
 C_n=A_0 n(n+1)+B_0,\qquad A_0,B_0\in\mathbb Q_{>0}.
\]
The same holds after an integer shift of the index.
\end{theorem}
\begin{proof}
Multiplying both states by a fixed positive integer clears the
coefficient denominators. Thus take $C_n=A n(n+1)+B$ with $A,B$
positive integers. The gcd chain now divides $2A$ and stabilizes at
$g$. As it divides every subsequent $C_n$ and $2A$, it also divides
$B$, since $A n(n+1)$ is an integer multiple of $2A$.
After division by $g$, the primitive numerator is
\[
 U_n=m\frac{n(n+1)}2+c,\qquad m=2A/g\in\mathbb Z_{>0},
 \quad c=B/g\in\mathbb Z_{>0}.
\]
Adjacent coprimality forces $c=1$: if $p\mid c$, choose arbitrarily
late $n$ divisible by $2p$; both $U_{n-1}$ and $U_n$ are then
divisible by $p$. If $m$ is odd, at $n\equiv1\pmod4$ the consecutive
values $U_n,U_{n+1}$ are both even. Hence $m$ is even and
\[
 U_n=k n(n+1)+1,\qquad k\in\mathbb Z_{>0}.
\]

For $f(X)=kX(X+1)+1$, put $\Delta=k(k-4)$ and let $\alpha$ be a
root. Direct substitution gives
\[
 -f(\alpha-1)f(\alpha+1)=\Delta-k^2=-4k.
\]
At every root modulo any good odd prime, the same expression must
be a square by equation~\eqref{eq:square}, using an arbitrarily late
index in that residue class. Lemma~\ref{lem:specialise} therefore
forces $-4k$ to be a square in $\mathbb Q(\alpha)$.
If $k\ge4$, this field is a subfield of $\mathbb R$ (including
$\mathbb Q$ in the repeated-root or rational-root cases), where the
negative number $-4k$ cannot be a square. If $k=1$, the field is
$\mathbb Q(\sqrt{-3})$, in which $-4$ is not a square. If $k=2$,
the field is $\mathbb Q(i)$, in which $-8$ is not a square.
For completeness, a rational number $h$ is a square in a quadratic
field $\mathbb Q(\sqrt d)$ only if $h$ or $h/d$ is a square in
$\mathbb Q$, by expanding $(u+v\sqrt d)^2$.

The sole remaining case is $k=3$. At $n\equiv1\pmod7$, the four
values $U_{n-2},U_{n-1},U_n,U_{n+1}$ are
\[
 (1,1,0,5)\pmod7.
\]
Thus $V_n\equiv2$ and equation~\eqref{eq:Didentity} gives
$V_{n-1}^2\equiv2$, so $V_{n-1}\equiv3$ or $4$.
The preceding step would require a root of
$x^2+x-V_{n-1}$ modulo $7$, whose discriminant is respectively
$6$ or $3$. Both are nonsquares modulo $7$. This eliminates $k=3$
and proves the theorem.
\end{proof}

\paragraph{What remains.}
The original near-squaring assumption yields subexponential integer
numerators and relative gaps tending to zero. It does not yield an
exact quadratic polynomial profile, nor a monic polynomial after
rational scaling. The second theorem also leaves negative offsets and
general rational linear coefficients outside its additional class.
No infinite counterexample has been constructed, and
no exclusion of all nonpolynomial orbits is proved here.
\end{document}
